\chapter{Come rifare ogni numero}\label{chapter:appendice-riproducibilita}
\sloppy

Ogni numero di questa tesi viene da un file dell'archivio \texttt{risultati/}, nel
repository della tesi, e ogni file registra il commit del codice che l'ha prodotto
e se l'albero era pulito (sezione~\ref{sec:principi}). Questa appendice dice come
si rifà un file partendo da zero.

\section*{I tre repository}

\begin{itemize}
  \item \textbf{Koskidex}: il motore, con \texttt{scripts/evaluate},
        \texttt{scripts/compare} e \texttt{scripts/pool}, e il diario delle
        modifiche al ranking in \texttt{eval/DIARIO.md}. Richiede Go e nessuna
        dipendenza oltre \texttt{golang.org/x/text}.
  \item \textbf{Documentale}, ramo \texttt{martin/tesi-magistrale}: i comandi di
        importazione, esportazione e valutazione aggiunti per la tesi, e
        l'innesto di Koskidex scelto da \texttt{SEARCH\_BACKEND}.
  \item \textbf{Il repository della tesi}: l'archivio \texttt{risultati/}, un
        README e uno script di analisi per ogni esperimento, gli strumenti di
        raccolta e di annotazione, questo testo.
\end{itemize}

\section*{L'ambiente}

La macchina delle misure è descritta in \texttt{risultati/macchina.md}: Apple M2,
Go 1.27.1, PHP 8.5.5. MySQL 8.4 ed Elasticsearch 9.1.0 girano in Docker,
Elasticsearch a nodo singolo con 1 GB di heap e la sicurezza spenta; gli embedding
vengono da Ollama con il modello \texttt{bge-m3}, impronta \texttt{790764642607}.
Gli strumenti scrivono nell'archivio solo se la variabile \texttt{TESI\_RISULTATI}
indica la cartella \texttt{risultati/}; se manca, lo dicono.

\section*{I dati}

Nessun dato sta nei repository: si riscaricano.
\begin{enumerate}
  \item \textbf{SciFact e NFCorpus}: \texttt{eval/corpora/fetch.sh} di Koskidex
        scarica gli archivi di BEIR e ne controlla l'impronta MD5.
  \item \textbf{Gli albi pretori}: \texttt{eval/corpora/fetch-albo.py} scarica il
        feed di Crispiano con i PDF e l'API del Friuli Venezia Giulia, ed estrae
        il testo dei PDF con \texttt{pdftotext}. Le fonti si aggiornano, e un
        nuovo scaricamento può dare atti in più: il corpus usato nelle misure ha
        la sua impronta SHA-256 registrata nei file dell'archivio.
  \item \textbf{Dentro Documentale}: \texttt{app:import-albo-corpus} importa gli
        atti con il modello dei documenti dell'app;
        \texttt{app:export-to-elastic-search} li indicizza in Elasticsearch, e lo
        script \texttt{strumenti/indicizza-koskidex.sh} fa lo stesso su Koskidex
        passando dall'app; \texttt{app:export-eval-corpus} esporta il corpus in
        formato BEIR per le valutazioni piatte.
\end{enumerate}

\section*{Le misure}

\begin{itemize}
  \item \textbf{Koskidex piatto}, collezioni pubbliche e albi:
        \texttt{go run ./scripts/evaluate} con le impostazioni da riga di comando
        (recupero, punteggio, analisi, vettori, fusione). Ogni file dell'archivio
        riporta nel blocco \texttt{config} le opzioni esatte.
  \item \textbf{Dall'app}: \texttt{app:eval-run-queries} esegue un file di query
        con la ricerca di produzione, su Elasticsearch o su Koskidex secondo
        \texttt{SEARCH\_BACKEND}; \texttt{scripts/evaluate -rankings} ne calcola
        le metriche con lo stesso codice di Koskidex piatto.
  \item \textbf{Gli esperimenti}: ogni cartella di
        \texttt{risultati/esperimenti/} ha un README con domanda, metodo,
        previsioni ed esito, e uno script \texttt{analizza.py} che legge i file
        dell'archivio e scrive un file di esito con il proprio commit. Il README
        cita i file di partenza per nome, e rifare l'esito vuol dire rilanciare lo
        script sullo stesso commit.
\end{itemize}

Il README di \texttt{risultati/} riporta i comandi esatti, i formati dei file e il
registro di ogni file tolto, rifatto o rinominato.
